% !TEX root = ../../Main.tex
\section{Future Work}
\label{Section:FutureWork}

The limitations above set the order of the future work. The most useful next step is the one that would most change the value of the results.

The first step is a true held-out evaluation. Training would end before the reporting period starts. The choice among candidates would use only data from before that point. The reported figures would then be an out-of-sample estimate. They would no longer describe the same window on which the models were chosen. Of all possible changes, this one would do the most to make the claims stronger. The framework already supports it. It needs a campaign run under this stricter protocol, and it needs no new features in the framework.

The second step follows from the first. The procedure used here splits the window only once. A walk-forward scheme with several folds would train each fold on the data before it and evaluate it on the data after it. This would give an out-of-sample record for the whole period, and not just for one year. The learned weights could also be carried from one fold into the next, instead of starting each fold from a new initialisation. This would test whether an agent adapts as the market changes. The framework supports this, but it was not used in this campaign.

The third step is to report how much the results vary between training runs. The path-robustness protocol used here changes the starting balance. This measures how sensitive a result is to the entry path, but not how sensitive it is to the training itself. The next step would report the distribution of outcomes across the whole pool of candidates, and where the published model falls within it. A reader could then see how much of a result comes from the method and how much comes from chance.

Apart from the evaluation, what the agents did suggests three further directions. The portfolio in \Cref{Section:Portfolio} gives the seven agents equal weights. Weighting them by risk, or by their correlation structure, should do better. The seven agents are almost independent of each other, which suggests that the gain would be real. The observation holds thirty-eight values from a single point in time. The natural next step is a recurrent or attention-based encoder that reads a window of past data directly. This was part of the original plan for this work. Last, the framework already runs the same strategy on a live account. A live deployment over a long period would test whether the execution assumptions hold when the orders are real. No backtest can test this.

One improvement to the framework itself should also be noted. The execution engine charges the spread correctly, by filling orders at the ask and at the bid. However, the exported trade record does not show this cost as a separate field. To recover it, the indirect calculation of \Cref{Equation:SpreadRecovery} is needed. If every trade recorded the spread cost next to its commission and swap, the cost decomposition in \Cref{Tables:Costs} could be read directly and would not have to be inferred. It could also be reported for each position, and not only in total.

The USD/JPY result needs a separate note. The current explanation of the failure on this pair is that the price window is monotonic, which means the price moves in one direction only. Such a window gives no gradient towards changing the size of the exposure. It has not been shown that this is the whole explanation. The pair should be studied again once the evaluation protocol above is in place.